% Optional short paragraph: % Optional short paragraph: % Optional short paragraph: % Optional short paragraph: \input{ai_workflow}
% Optional appendix: \begingroup\def\StaffordWorkflowAppendix{}
% \input{ai_workflow}\endgroup
% The manuscript contains an inline adaptation for direct editing.

\ifdefined\StaffordWorkflowAppendix
% Evidence: project ledgers, archived outward session messages, model-use
% records, independent reviews, and canonical formal-proof provenance.
\section*{Appendix: AI-assisted research and formalization}

Under Christopher Albert's direction, AI agents supported proof discovery,
counterexample searches, review, and formalization. The work began with
cyclic-extension computations in Macaulay2 and Maple. Attempts to construct
universal cofactors directly produced special cases and reductions but
eventually stalled. The decisive change was to study the canonical quotient
$Q=A_n/(dA_n+x^N dA_n)$: Euler surjectivity and monicity constrained its
characteristic support, while Gabber's involutivity and asymptotic conormal
geometry forced that support to be empty. This turned cofactor selection
into a quotient-vanishing problem. Full Lean formalization was deferred
until the paper argument was complete, although small formal checks
continued throughout. In the final formal development, support avoidance
over the coordinate hyperplane was proved by a localized filtered Koszul
length contradiction.

The workflow separated proposed ideas from verified claims. One controller
maintained the authoritative proof state, while parallel workers explored
distinct mechanisms and returned evidence with an explicit scope.
Independent reviewers checked fixed revisions for quantifiers, sidedness,
circularity, and missing hypotheses. The project protocols followed the
same discipline as the shared \texttt{math-frontier} skill: freeze the
claim, identify its first unresolved dependency, and specify how a route
could fail. Conceptual prompts requested a small number of distinct
mechanisms; paper-only tasks sought complete arguments before implementation.
A condensed representative prompt, adapted from the archived protocols, is:
\begin{quote}\small
Work on one precisely stated unresolved claim. Propose at most three
distinct mechanisms, each with an intermediate lemma, a noncircular
justification, and an adversarial test. Preserve the quantifiers and
right-sided cofactors. For a paper task, supply a complete argument; for
formalization, prove the required hypotheses without adding axioms or
using \texttt{sorry}. Return evidence and the first unsupported bridge; do not
promote a conditional result to a theorem.
\end{quote}

Recorded research and review sessions used GPT-5.5, GPT-5.6 Luna and Sol,
GPT-6 Astra, and Claude Fable 5, Haiku 4.5 and Opus 5. The project protocols
assigned routine implementation and replay to Luna and conceptual
consultation to Sol; independent reports also record Opus reviews.
Subsequent manuscript, proof-map, build and release work used Fable 5.1,
Opus 5, Sonnet 5 and GPT-6 Luna. These records establish participation and task roles,
without assigning a particular mathematical idea to one model.

The completed proof uses Lean~4 and Mathlib and pins AlgebraicAnalysis v0.3.0
\cite{Lean4,Mathlib2020,AlgebraicAnalysis030}; the later library archive is v0.3.2~\cite{AlgebraicAnalysis032}.
Formalization distinguished theorems proving required inputs from
conditional theorems merely consuming them. Statement comparison and
independent NanoDa and Lean kernel checks verified the submitted declarations
\cite{Comparator,NanoDa,PalomarDocs}; their proofs use only
\texttt{propext}, \texttt{Classical.choice}, and \texttt{Quot.sound}.
The Stafford38 and Global Stafford headline proofs are registered separately in
Palomar~\cite{Stafford38Palomar2026,GlobalStaffordPalomar2026}; their Lean developments are archived~\cite{Stafford38Zenodo120,GlobalStaffordZenodo106}.
These checks certify the formal declarations; human review of the paper's
exposition and its correspondence with them remains a separate task.

\else
\paragraph{AI-assisted research and formalization.}
Under Christopher Albert's direction, AI agents supported discovery,
counterexample searches, independent review, and formalization. Early
cofactor-selection approaches stalled; the proof advanced by reframing
the problem as vanishing of the canonical quotient, using Euler
surjectivity, monicity and characteristic-support geometry. Full Lean
formalization followed the completed paper argument, while bounded checks
continued during discovery. Parallel workers returned scoped evidence to
one controller, and reviewers checked quantifiers, sidedness and unsupported
bridges independently. The Lean/Mathlib development uses AlgebraicAnalysis
\cite{Lean4,Mathlib2020,AlgebraicAnalysis030}; its registered proof and
statement/kernel checks are recorded by Palomar
\cite{Stafford38Palomar2026,Comparator,NanoDa,PalomarDocs}.
The appendix describes the models, proof history, and prompt strategy.
\fi

% Optional appendix: \begingroup\def\StaffordWorkflowAppendix{}
% % Optional short paragraph: \input{ai_workflow}
% Optional appendix: \begingroup\def\StaffordWorkflowAppendix{}
% \input{ai_workflow}\endgroup
% The manuscript contains an inline adaptation for direct editing.

\ifdefined\StaffordWorkflowAppendix
% Evidence: project ledgers, archived outward session messages, model-use
% records, independent reviews, and canonical formal-proof provenance.
\section*{Appendix: AI-assisted research and formalization}

Under Christopher Albert's direction, AI agents supported proof discovery,
counterexample searches, review, and formalization. The work began with
cyclic-extension computations in Macaulay2 and Maple. Attempts to construct
universal cofactors directly produced special cases and reductions but
eventually stalled. The decisive change was to study the canonical quotient
$Q=A_n/(dA_n+x^N dA_n)$: Euler surjectivity and monicity constrained its
characteristic support, while Gabber's involutivity and asymptotic conormal
geometry forced that support to be empty. This turned cofactor selection
into a quotient-vanishing problem. Full Lean formalization was deferred
until the paper argument was complete, although small formal checks
continued throughout. In the final formal development, support avoidance
over the coordinate hyperplane was proved by a localized filtered Koszul
length contradiction.

The workflow separated proposed ideas from verified claims. One controller
maintained the authoritative proof state, while parallel workers explored
distinct mechanisms and returned evidence with an explicit scope.
Independent reviewers checked fixed revisions for quantifiers, sidedness,
circularity, and missing hypotheses. The project protocols followed the
same discipline as the shared \texttt{math-frontier} skill: freeze the
claim, identify its first unresolved dependency, and specify how a route
could fail. Conceptual prompts requested a small number of distinct
mechanisms; paper-only tasks sought complete arguments before implementation.
A condensed representative prompt, adapted from the archived protocols, is:
\begin{quote}\small
Work on one precisely stated unresolved claim. Propose at most three
distinct mechanisms, each with an intermediate lemma, a noncircular
justification, and an adversarial test. Preserve the quantifiers and
right-sided cofactors. For a paper task, supply a complete argument; for
formalization, prove the required hypotheses without adding axioms or
using \texttt{sorry}. Return evidence and the first unsupported bridge; do not
promote a conditional result to a theorem.
\end{quote}

Recorded research and review sessions used GPT-5.5, GPT-5.6 Luna and Sol,
GPT-6 Astra, and Claude Fable 5, Haiku 4.5 and Opus 5. The project protocols
assigned routine implementation and replay to Luna and conceptual
consultation to Sol; independent reports also record Opus reviews.
Subsequent manuscript, proof-map, build and release work used Fable 5.1,
Opus 5, Sonnet 5 and GPT-6 Luna. These records establish participation and task roles,
without assigning a particular mathematical idea to one model.

The completed proof uses Lean~4 and Mathlib and pins AlgebraicAnalysis v0.3.0
\cite{Lean4,Mathlib2020,AlgebraicAnalysis030}; the later library archive is v0.3.2~\cite{AlgebraicAnalysis032}.
Formalization distinguished theorems proving required inputs from
conditional theorems merely consuming them. Statement comparison and
independent NanoDa and Lean kernel checks verified the submitted declarations
\cite{Comparator,NanoDa,PalomarDocs}; their proofs use only
\texttt{propext}, \texttt{Classical.choice}, and \texttt{Quot.sound}.
The Stafford38 and Global Stafford headline proofs are registered separately in
Palomar~\cite{Stafford38Palomar2026,GlobalStaffordPalomar2026}; their Lean developments are archived~\cite{Stafford38Zenodo120,GlobalStaffordZenodo106}.
These checks certify the formal declarations; human review of the paper's
exposition and its correspondence with them remains a separate task.

\else
\paragraph{AI-assisted research and formalization.}
Under Christopher Albert's direction, AI agents supported discovery,
counterexample searches, independent review, and formalization. Early
cofactor-selection approaches stalled; the proof advanced by reframing
the problem as vanishing of the canonical quotient, using Euler
surjectivity, monicity and characteristic-support geometry. Full Lean
formalization followed the completed paper argument, while bounded checks
continued during discovery. Parallel workers returned scoped evidence to
one controller, and reviewers checked quantifiers, sidedness and unsupported
bridges independently. The Lean/Mathlib development uses AlgebraicAnalysis
\cite{Lean4,Mathlib2020,AlgebraicAnalysis030}; its registered proof and
statement/kernel checks are recorded by Palomar
\cite{Stafford38Palomar2026,Comparator,NanoDa,PalomarDocs}.
The appendix describes the models, proof history, and prompt strategy.
\fi
\endgroup
% The manuscript contains an inline adaptation for direct editing.

\ifdefined\StaffordWorkflowAppendix
% Evidence: project ledgers, archived outward session messages, model-use
% records, independent reviews, and canonical formal-proof provenance.
\section*{Appendix: AI-assisted research and formalization}

Under Christopher Albert's direction, AI agents supported proof discovery,
counterexample searches, review, and formalization. The work began with
cyclic-extension computations in Macaulay2 and Maple. Attempts to construct
universal cofactors directly produced special cases and reductions but
eventually stalled. The decisive change was to study the canonical quotient
$Q=A_n/(dA_n+x^N dA_n)$: Euler surjectivity and monicity constrained its
characteristic support, while Gabber's involutivity and asymptotic conormal
geometry forced that support to be empty. This turned cofactor selection
into a quotient-vanishing problem. Full Lean formalization was deferred
until the paper argument was complete, although small formal checks
continued throughout. In the final formal development, support avoidance
over the coordinate hyperplane was proved by a localized filtered Koszul
length contradiction.

The workflow separated proposed ideas from verified claims. One controller
maintained the authoritative proof state, while parallel workers explored
distinct mechanisms and returned evidence with an explicit scope.
Independent reviewers checked fixed revisions for quantifiers, sidedness,
circularity, and missing hypotheses. The project protocols followed the
same discipline as the shared \texttt{math-frontier} skill: freeze the
claim, identify its first unresolved dependency, and specify how a route
could fail. Conceptual prompts requested a small number of distinct
mechanisms; paper-only tasks sought complete arguments before implementation.
A condensed representative prompt, adapted from the archived protocols, is:
\begin{quote}\small
Work on one precisely stated unresolved claim. Propose at most three
distinct mechanisms, each with an intermediate lemma, a noncircular
justification, and an adversarial test. Preserve the quantifiers and
right-sided cofactors. For a paper task, supply a complete argument; for
formalization, prove the required hypotheses without adding axioms or
using \texttt{sorry}. Return evidence and the first unsupported bridge; do not
promote a conditional result to a theorem.
\end{quote}

Recorded research and review sessions used GPT-5.5, GPT-5.6 Luna and Sol,
GPT-6 Astra, and Claude Fable 5, Haiku 4.5 and Opus 5. The project protocols
assigned routine implementation and replay to Luna and conceptual
consultation to Sol; independent reports also record Opus reviews.
Subsequent manuscript, proof-map, build and release work used Fable 5.1,
Opus 5, Sonnet 5 and GPT-6 Luna. These records establish participation and task roles,
without assigning a particular mathematical idea to one model.

The completed proof uses Lean~4 and Mathlib and pins AlgebraicAnalysis v0.3.0
\cite{Lean4,Mathlib2020,AlgebraicAnalysis030}; the later library archive is v0.3.2~\cite{AlgebraicAnalysis032}.
Formalization distinguished theorems proving required inputs from
conditional theorems merely consuming them. Statement comparison and
independent NanoDa and Lean kernel checks verified the submitted declarations
\cite{Comparator,NanoDa,PalomarDocs}; their proofs use only
\texttt{propext}, \texttt{Classical.choice}, and \texttt{Quot.sound}.
The Stafford38 and Global Stafford headline proofs are registered separately in
Palomar~\cite{Stafford38Palomar2026,GlobalStaffordPalomar2026}; their Lean developments are archived~\cite{Stafford38Zenodo120,GlobalStaffordZenodo106}.
These checks certify the formal declarations; human review of the paper's
exposition and its correspondence with them remains a separate task.

\else
\paragraph{AI-assisted research and formalization.}
Under Christopher Albert's direction, AI agents supported discovery,
counterexample searches, independent review, and formalization. Early
cofactor-selection approaches stalled; the proof advanced by reframing
the problem as vanishing of the canonical quotient, using Euler
surjectivity, monicity and characteristic-support geometry. Full Lean
formalization followed the completed paper argument, while bounded checks
continued during discovery. Parallel workers returned scoped evidence to
one controller, and reviewers checked quantifiers, sidedness and unsupported
bridges independently. The Lean/Mathlib development uses AlgebraicAnalysis
\cite{Lean4,Mathlib2020,AlgebraicAnalysis030}; its registered proof and
statement/kernel checks are recorded by Palomar
\cite{Stafford38Palomar2026,Comparator,NanoDa,PalomarDocs}.
The appendix describes the models, proof history, and prompt strategy.
\fi

% Optional appendix: \begingroup\def\StaffordWorkflowAppendix{}
% % Optional short paragraph: % Optional short paragraph: \input{ai_workflow}
% Optional appendix: \begingroup\def\StaffordWorkflowAppendix{}
% \input{ai_workflow}\endgroup
% The manuscript contains an inline adaptation for direct editing.

\ifdefined\StaffordWorkflowAppendix
% Evidence: project ledgers, archived outward session messages, model-use
% records, independent reviews, and canonical formal-proof provenance.
\section*{Appendix: AI-assisted research and formalization}

Under Christopher Albert's direction, AI agents supported proof discovery,
counterexample searches, review, and formalization. The work began with
cyclic-extension computations in Macaulay2 and Maple. Attempts to construct
universal cofactors directly produced special cases and reductions but
eventually stalled. The decisive change was to study the canonical quotient
$Q=A_n/(dA_n+x^N dA_n)$: Euler surjectivity and monicity constrained its
characteristic support, while Gabber's involutivity and asymptotic conormal
geometry forced that support to be empty. This turned cofactor selection
into a quotient-vanishing problem. Full Lean formalization was deferred
until the paper argument was complete, although small formal checks
continued throughout. In the final formal development, support avoidance
over the coordinate hyperplane was proved by a localized filtered Koszul
length contradiction.

The workflow separated proposed ideas from verified claims. One controller
maintained the authoritative proof state, while parallel workers explored
distinct mechanisms and returned evidence with an explicit scope.
Independent reviewers checked fixed revisions for quantifiers, sidedness,
circularity, and missing hypotheses. The project protocols followed the
same discipline as the shared \texttt{math-frontier} skill: freeze the
claim, identify its first unresolved dependency, and specify how a route
could fail. Conceptual prompts requested a small number of distinct
mechanisms; paper-only tasks sought complete arguments before implementation.
A condensed representative prompt, adapted from the archived protocols, is:
\begin{quote}\small
Work on one precisely stated unresolved claim. Propose at most three
distinct mechanisms, each with an intermediate lemma, a noncircular
justification, and an adversarial test. Preserve the quantifiers and
right-sided cofactors. For a paper task, supply a complete argument; for
formalization, prove the required hypotheses without adding axioms or
using \texttt{sorry}. Return evidence and the first unsupported bridge; do not
promote a conditional result to a theorem.
\end{quote}

Recorded research and review sessions used GPT-5.5, GPT-5.6 Luna and Sol,
GPT-6 Astra, and Claude Fable 5, Haiku 4.5 and Opus 5. The project protocols
assigned routine implementation and replay to Luna and conceptual
consultation to Sol; independent reports also record Opus reviews.
Subsequent manuscript, proof-map, build and release work used Fable 5.1,
Opus 5, Sonnet 5 and GPT-6 Luna. These records establish participation and task roles,
without assigning a particular mathematical idea to one model.

The completed proof uses Lean~4 and Mathlib and pins AlgebraicAnalysis v0.3.0
\cite{Lean4,Mathlib2020,AlgebraicAnalysis030}; the later library archive is v0.3.2~\cite{AlgebraicAnalysis032}.
Formalization distinguished theorems proving required inputs from
conditional theorems merely consuming them. Statement comparison and
independent NanoDa and Lean kernel checks verified the submitted declarations
\cite{Comparator,NanoDa,PalomarDocs}; their proofs use only
\texttt{propext}, \texttt{Classical.choice}, and \texttt{Quot.sound}.
The Stafford38 and Global Stafford headline proofs are registered separately in
Palomar~\cite{Stafford38Palomar2026,GlobalStaffordPalomar2026}; their Lean developments are archived~\cite{Stafford38Zenodo120,GlobalStaffordZenodo106}.
These checks certify the formal declarations; human review of the paper's
exposition and its correspondence with them remains a separate task.

\else
\paragraph{AI-assisted research and formalization.}
Under Christopher Albert's direction, AI agents supported discovery,
counterexample searches, independent review, and formalization. Early
cofactor-selection approaches stalled; the proof advanced by reframing
the problem as vanishing of the canonical quotient, using Euler
surjectivity, monicity and characteristic-support geometry. Full Lean
formalization followed the completed paper argument, while bounded checks
continued during discovery. Parallel workers returned scoped evidence to
one controller, and reviewers checked quantifiers, sidedness and unsupported
bridges independently. The Lean/Mathlib development uses AlgebraicAnalysis
\cite{Lean4,Mathlib2020,AlgebraicAnalysis030}; its registered proof and
statement/kernel checks are recorded by Palomar
\cite{Stafford38Palomar2026,Comparator,NanoDa,PalomarDocs}.
The appendix describes the models, proof history, and prompt strategy.
\fi

% Optional appendix: \begingroup\def\StaffordWorkflowAppendix{}
% % Optional short paragraph: \input{ai_workflow}
% Optional appendix: \begingroup\def\StaffordWorkflowAppendix{}
% \input{ai_workflow}\endgroup
% The manuscript contains an inline adaptation for direct editing.

\ifdefined\StaffordWorkflowAppendix
% Evidence: project ledgers, archived outward session messages, model-use
% records, independent reviews, and canonical formal-proof provenance.
\section*{Appendix: AI-assisted research and formalization}

Under Christopher Albert's direction, AI agents supported proof discovery,
counterexample searches, review, and formalization. The work began with
cyclic-extension computations in Macaulay2 and Maple. Attempts to construct
universal cofactors directly produced special cases and reductions but
eventually stalled. The decisive change was to study the canonical quotient
$Q=A_n/(dA_n+x^N dA_n)$: Euler surjectivity and monicity constrained its
characteristic support, while Gabber's involutivity and asymptotic conormal
geometry forced that support to be empty. This turned cofactor selection
into a quotient-vanishing problem. Full Lean formalization was deferred
until the paper argument was complete, although small formal checks
continued throughout. In the final formal development, support avoidance
over the coordinate hyperplane was proved by a localized filtered Koszul
length contradiction.

The workflow separated proposed ideas from verified claims. One controller
maintained the authoritative proof state, while parallel workers explored
distinct mechanisms and returned evidence with an explicit scope.
Independent reviewers checked fixed revisions for quantifiers, sidedness,
circularity, and missing hypotheses. The project protocols followed the
same discipline as the shared \texttt{math-frontier} skill: freeze the
claim, identify its first unresolved dependency, and specify how a route
could fail. Conceptual prompts requested a small number of distinct
mechanisms; paper-only tasks sought complete arguments before implementation.
A condensed representative prompt, adapted from the archived protocols, is:
\begin{quote}\small
Work on one precisely stated unresolved claim. Propose at most three
distinct mechanisms, each with an intermediate lemma, a noncircular
justification, and an adversarial test. Preserve the quantifiers and
right-sided cofactors. For a paper task, supply a complete argument; for
formalization, prove the required hypotheses without adding axioms or
using \texttt{sorry}. Return evidence and the first unsupported bridge; do not
promote a conditional result to a theorem.
\end{quote}

Recorded research and review sessions used GPT-5.5, GPT-5.6 Luna and Sol,
GPT-6 Astra, and Claude Fable 5, Haiku 4.5 and Opus 5. The project protocols
assigned routine implementation and replay to Luna and conceptual
consultation to Sol; independent reports also record Opus reviews.
Subsequent manuscript, proof-map, build and release work used Fable 5.1,
Opus 5, Sonnet 5 and GPT-6 Luna. These records establish participation and task roles,
without assigning a particular mathematical idea to one model.

The completed proof uses Lean~4 and Mathlib and pins AlgebraicAnalysis v0.3.0
\cite{Lean4,Mathlib2020,AlgebraicAnalysis030}; the later library archive is v0.3.2~\cite{AlgebraicAnalysis032}.
Formalization distinguished theorems proving required inputs from
conditional theorems merely consuming them. Statement comparison and
independent NanoDa and Lean kernel checks verified the submitted declarations
\cite{Comparator,NanoDa,PalomarDocs}; their proofs use only
\texttt{propext}, \texttt{Classical.choice}, and \texttt{Quot.sound}.
The Stafford38 and Global Stafford headline proofs are registered separately in
Palomar~\cite{Stafford38Palomar2026,GlobalStaffordPalomar2026}; their Lean developments are archived~\cite{Stafford38Zenodo120,GlobalStaffordZenodo106}.
These checks certify the formal declarations; human review of the paper's
exposition and its correspondence with them remains a separate task.

\else
\paragraph{AI-assisted research and formalization.}
Under Christopher Albert's direction, AI agents supported discovery,
counterexample searches, independent review, and formalization. Early
cofactor-selection approaches stalled; the proof advanced by reframing
the problem as vanishing of the canonical quotient, using Euler
surjectivity, monicity and characteristic-support geometry. Full Lean
formalization followed the completed paper argument, while bounded checks
continued during discovery. Parallel workers returned scoped evidence to
one controller, and reviewers checked quantifiers, sidedness and unsupported
bridges independently. The Lean/Mathlib development uses AlgebraicAnalysis
\cite{Lean4,Mathlib2020,AlgebraicAnalysis030}; its registered proof and
statement/kernel checks are recorded by Palomar
\cite{Stafford38Palomar2026,Comparator,NanoDa,PalomarDocs}.
The appendix describes the models, proof history, and prompt strategy.
\fi
\endgroup
% The manuscript contains an inline adaptation for direct editing.

\ifdefined\StaffordWorkflowAppendix
% Evidence: project ledgers, archived outward session messages, model-use
% records, independent reviews, and canonical formal-proof provenance.
\section*{Appendix: AI-assisted research and formalization}

Under Christopher Albert's direction, AI agents supported proof discovery,
counterexample searches, review, and formalization. The work began with
cyclic-extension computations in Macaulay2 and Maple. Attempts to construct
universal cofactors directly produced special cases and reductions but
eventually stalled. The decisive change was to study the canonical quotient
$Q=A_n/(dA_n+x^N dA_n)$: Euler surjectivity and monicity constrained its
characteristic support, while Gabber's involutivity and asymptotic conormal
geometry forced that support to be empty. This turned cofactor selection
into a quotient-vanishing problem. Full Lean formalization was deferred
until the paper argument was complete, although small formal checks
continued throughout. In the final formal development, support avoidance
over the coordinate hyperplane was proved by a localized filtered Koszul
length contradiction.

The workflow separated proposed ideas from verified claims. One controller
maintained the authoritative proof state, while parallel workers explored
distinct mechanisms and returned evidence with an explicit scope.
Independent reviewers checked fixed revisions for quantifiers, sidedness,
circularity, and missing hypotheses. The project protocols followed the
same discipline as the shared \texttt{math-frontier} skill: freeze the
claim, identify its first unresolved dependency, and specify how a route
could fail. Conceptual prompts requested a small number of distinct
mechanisms; paper-only tasks sought complete arguments before implementation.
A condensed representative prompt, adapted from the archived protocols, is:
\begin{quote}\small
Work on one precisely stated unresolved claim. Propose at most three
distinct mechanisms, each with an intermediate lemma, a noncircular
justification, and an adversarial test. Preserve the quantifiers and
right-sided cofactors. For a paper task, supply a complete argument; for
formalization, prove the required hypotheses without adding axioms or
using \texttt{sorry}. Return evidence and the first unsupported bridge; do not
promote a conditional result to a theorem.
\end{quote}

Recorded research and review sessions used GPT-5.5, GPT-5.6 Luna and Sol,
GPT-6 Astra, and Claude Fable 5, Haiku 4.5 and Opus 5. The project protocols
assigned routine implementation and replay to Luna and conceptual
consultation to Sol; independent reports also record Opus reviews.
Subsequent manuscript, proof-map, build and release work used Fable 5.1,
Opus 5, Sonnet 5 and GPT-6 Luna. These records establish participation and task roles,
without assigning a particular mathematical idea to one model.

The completed proof uses Lean~4 and Mathlib and pins AlgebraicAnalysis v0.3.0
\cite{Lean4,Mathlib2020,AlgebraicAnalysis030}; the later library archive is v0.3.2~\cite{AlgebraicAnalysis032}.
Formalization distinguished theorems proving required inputs from
conditional theorems merely consuming them. Statement comparison and
independent NanoDa and Lean kernel checks verified the submitted declarations
\cite{Comparator,NanoDa,PalomarDocs}; their proofs use only
\texttt{propext}, \texttt{Classical.choice}, and \texttt{Quot.sound}.
The Stafford38 and Global Stafford headline proofs are registered separately in
Palomar~\cite{Stafford38Palomar2026,GlobalStaffordPalomar2026}; their Lean developments are archived~\cite{Stafford38Zenodo120,GlobalStaffordZenodo106}.
These checks certify the formal declarations; human review of the paper's
exposition and its correspondence with them remains a separate task.

\else
\paragraph{AI-assisted research and formalization.}
Under Christopher Albert's direction, AI agents supported discovery,
counterexample searches, independent review, and formalization. Early
cofactor-selection approaches stalled; the proof advanced by reframing
the problem as vanishing of the canonical quotient, using Euler
surjectivity, monicity and characteristic-support geometry. Full Lean
formalization followed the completed paper argument, while bounded checks
continued during discovery. Parallel workers returned scoped evidence to
one controller, and reviewers checked quantifiers, sidedness and unsupported
bridges independently. The Lean/Mathlib development uses AlgebraicAnalysis
\cite{Lean4,Mathlib2020,AlgebraicAnalysis030}; its registered proof and
statement/kernel checks are recorded by Palomar
\cite{Stafford38Palomar2026,Comparator,NanoDa,PalomarDocs}.
The appendix describes the models, proof history, and prompt strategy.
\fi
\endgroup
% The manuscript contains an inline adaptation for direct editing.

\ifdefined\StaffordWorkflowAppendix
% Evidence: project ledgers, archived outward session messages, model-use
% records, independent reviews, and canonical formal-proof provenance.
\section*{Appendix: AI-assisted research and formalization}

Under Christopher Albert's direction, AI agents supported proof discovery,
counterexample searches, review, and formalization. The work began with
cyclic-extension computations in Macaulay2 and Maple. Attempts to construct
universal cofactors directly produced special cases and reductions but
eventually stalled. The decisive change was to study the canonical quotient
$Q=A_n/(dA_n+x^N dA_n)$: Euler surjectivity and monicity constrained its
characteristic support, while Gabber's involutivity and asymptotic conormal
geometry forced that support to be empty. This turned cofactor selection
into a quotient-vanishing problem. Full Lean formalization was deferred
until the paper argument was complete, although small formal checks
continued throughout. In the final formal development, support avoidance
over the coordinate hyperplane was proved by a localized filtered Koszul
length contradiction.

The workflow separated proposed ideas from verified claims. One controller
maintained the authoritative proof state, while parallel workers explored
distinct mechanisms and returned evidence with an explicit scope.
Independent reviewers checked fixed revisions for quantifiers, sidedness,
circularity, and missing hypotheses. The project protocols followed the
same discipline as the shared \texttt{math-frontier} skill: freeze the
claim, identify its first unresolved dependency, and specify how a route
could fail. Conceptual prompts requested a small number of distinct
mechanisms; paper-only tasks sought complete arguments before implementation.
A condensed representative prompt, adapted from the archived protocols, is:
\begin{quote}\small
Work on one precisely stated unresolved claim. Propose at most three
distinct mechanisms, each with an intermediate lemma, a noncircular
justification, and an adversarial test. Preserve the quantifiers and
right-sided cofactors. For a paper task, supply a complete argument; for
formalization, prove the required hypotheses without adding axioms or
using \texttt{sorry}. Return evidence and the first unsupported bridge; do not
promote a conditional result to a theorem.
\end{quote}

Recorded research and review sessions used GPT-5.5, GPT-5.6 Luna and Sol,
GPT-6 Astra, and Claude Fable 5, Haiku 4.5 and Opus 5. The project protocols
assigned routine implementation and replay to Luna and conceptual
consultation to Sol; independent reports also record Opus reviews.
Subsequent manuscript, proof-map, build and release work used Fable 5.1,
Opus 5, Sonnet 5 and GPT-6 Luna. These records establish participation and task roles,
without assigning a particular mathematical idea to one model.

The completed proof uses Lean~4 and Mathlib and pins AlgebraicAnalysis v0.3.0
\cite{Lean4,Mathlib2020,AlgebraicAnalysis030}; the later library archive is v0.3.2~\cite{AlgebraicAnalysis032}.
Formalization distinguished theorems proving required inputs from
conditional theorems merely consuming them. Statement comparison and
independent NanoDa and Lean kernel checks verified the submitted declarations
\cite{Comparator,NanoDa,PalomarDocs}; their proofs use only
\texttt{propext}, \texttt{Classical.choice}, and \texttt{Quot.sound}.
The Stafford38 and Global Stafford headline proofs are registered separately in
Palomar~\cite{Stafford38Palomar2026,GlobalStaffordPalomar2026}; their Lean developments are archived~\cite{Stafford38Zenodo120,GlobalStaffordZenodo106}.
These checks certify the formal declarations; human review of the paper's
exposition and its correspondence with them remains a separate task.

\else
\paragraph{AI-assisted research and formalization.}
Under Christopher Albert's direction, AI agents supported discovery,
counterexample searches, independent review, and formalization. Early
cofactor-selection approaches stalled; the proof advanced by reframing
the problem as vanishing of the canonical quotient, using Euler
surjectivity, monicity and characteristic-support geometry. Full Lean
formalization followed the completed paper argument, while bounded checks
continued during discovery. Parallel workers returned scoped evidence to
one controller, and reviewers checked quantifiers, sidedness and unsupported
bridges independently. The Lean/Mathlib development uses AlgebraicAnalysis
\cite{Lean4,Mathlib2020,AlgebraicAnalysis030}; its registered proof and
statement/kernel checks are recorded by Palomar
\cite{Stafford38Palomar2026,Comparator,NanoDa,PalomarDocs}.
The appendix describes the models, proof history, and prompt strategy.
\fi

% Optional appendix: \begingroup\def\StaffordWorkflowAppendix{}
% % Optional short paragraph: % Optional short paragraph: % Optional short paragraph: \input{ai_workflow}
% Optional appendix: \begingroup\def\StaffordWorkflowAppendix{}
% \input{ai_workflow}\endgroup
% The manuscript contains an inline adaptation for direct editing.

\ifdefined\StaffordWorkflowAppendix
% Evidence: project ledgers, archived outward session messages, model-use
% records, independent reviews, and canonical formal-proof provenance.
\section*{Appendix: AI-assisted research and formalization}

Under Christopher Albert's direction, AI agents supported proof discovery,
counterexample searches, review, and formalization. The work began with
cyclic-extension computations in Macaulay2 and Maple. Attempts to construct
universal cofactors directly produced special cases and reductions but
eventually stalled. The decisive change was to study the canonical quotient
$Q=A_n/(dA_n+x^N dA_n)$: Euler surjectivity and monicity constrained its
characteristic support, while Gabber's involutivity and asymptotic conormal
geometry forced that support to be empty. This turned cofactor selection
into a quotient-vanishing problem. Full Lean formalization was deferred
until the paper argument was complete, although small formal checks
continued throughout. In the final formal development, support avoidance
over the coordinate hyperplane was proved by a localized filtered Koszul
length contradiction.

The workflow separated proposed ideas from verified claims. One controller
maintained the authoritative proof state, while parallel workers explored
distinct mechanisms and returned evidence with an explicit scope.
Independent reviewers checked fixed revisions for quantifiers, sidedness,
circularity, and missing hypotheses. The project protocols followed the
same discipline as the shared \texttt{math-frontier} skill: freeze the
claim, identify its first unresolved dependency, and specify how a route
could fail. Conceptual prompts requested a small number of distinct
mechanisms; paper-only tasks sought complete arguments before implementation.
A condensed representative prompt, adapted from the archived protocols, is:
\begin{quote}\small
Work on one precisely stated unresolved claim. Propose at most three
distinct mechanisms, each with an intermediate lemma, a noncircular
justification, and an adversarial test. Preserve the quantifiers and
right-sided cofactors. For a paper task, supply a complete argument; for
formalization, prove the required hypotheses without adding axioms or
using \texttt{sorry}. Return evidence and the first unsupported bridge; do not
promote a conditional result to a theorem.
\end{quote}

Recorded research and review sessions used GPT-5.5, GPT-5.6 Luna and Sol,
GPT-6 Astra, and Claude Fable 5, Haiku 4.5 and Opus 5. The project protocols
assigned routine implementation and replay to Luna and conceptual
consultation to Sol; independent reports also record Opus reviews.
Subsequent manuscript, proof-map, build and release work used Fable 5.1,
Opus 5, Sonnet 5 and GPT-6 Luna. These records establish participation and task roles,
without assigning a particular mathematical idea to one model.

The completed proof uses Lean~4 and Mathlib and pins AlgebraicAnalysis v0.3.0
\cite{Lean4,Mathlib2020,AlgebraicAnalysis030}; the later library archive is v0.3.2~\cite{AlgebraicAnalysis032}.
Formalization distinguished theorems proving required inputs from
conditional theorems merely consuming them. Statement comparison and
independent NanoDa and Lean kernel checks verified the submitted declarations
\cite{Comparator,NanoDa,PalomarDocs}; their proofs use only
\texttt{propext}, \texttt{Classical.choice}, and \texttt{Quot.sound}.
The Stafford38 and Global Stafford headline proofs are registered separately in
Palomar~\cite{Stafford38Palomar2026,GlobalStaffordPalomar2026}; their Lean developments are archived~\cite{Stafford38Zenodo120,GlobalStaffordZenodo106}.
These checks certify the formal declarations; human review of the paper's
exposition and its correspondence with them remains a separate task.

\else
\paragraph{AI-assisted research and formalization.}
Under Christopher Albert's direction, AI agents supported discovery,
counterexample searches, independent review, and formalization. Early
cofactor-selection approaches stalled; the proof advanced by reframing
the problem as vanishing of the canonical quotient, using Euler
surjectivity, monicity and characteristic-support geometry. Full Lean
formalization followed the completed paper argument, while bounded checks
continued during discovery. Parallel workers returned scoped evidence to
one controller, and reviewers checked quantifiers, sidedness and unsupported
bridges independently. The Lean/Mathlib development uses AlgebraicAnalysis
\cite{Lean4,Mathlib2020,AlgebraicAnalysis030}; its registered proof and
statement/kernel checks are recorded by Palomar
\cite{Stafford38Palomar2026,Comparator,NanoDa,PalomarDocs}.
The appendix describes the models, proof history, and prompt strategy.
\fi

% Optional appendix: \begingroup\def\StaffordWorkflowAppendix{}
% % Optional short paragraph: \input{ai_workflow}
% Optional appendix: \begingroup\def\StaffordWorkflowAppendix{}
% \input{ai_workflow}\endgroup
% The manuscript contains an inline adaptation for direct editing.

\ifdefined\StaffordWorkflowAppendix
% Evidence: project ledgers, archived outward session messages, model-use
% records, independent reviews, and canonical formal-proof provenance.
\section*{Appendix: AI-assisted research and formalization}

Under Christopher Albert's direction, AI agents supported proof discovery,
counterexample searches, review, and formalization. The work began with
cyclic-extension computations in Macaulay2 and Maple. Attempts to construct
universal cofactors directly produced special cases and reductions but
eventually stalled. The decisive change was to study the canonical quotient
$Q=A_n/(dA_n+x^N dA_n)$: Euler surjectivity and monicity constrained its
characteristic support, while Gabber's involutivity and asymptotic conormal
geometry forced that support to be empty. This turned cofactor selection
into a quotient-vanishing problem. Full Lean formalization was deferred
until the paper argument was complete, although small formal checks
continued throughout. In the final formal development, support avoidance
over the coordinate hyperplane was proved by a localized filtered Koszul
length contradiction.

The workflow separated proposed ideas from verified claims. One controller
maintained the authoritative proof state, while parallel workers explored
distinct mechanisms and returned evidence with an explicit scope.
Independent reviewers checked fixed revisions for quantifiers, sidedness,
circularity, and missing hypotheses. The project protocols followed the
same discipline as the shared \texttt{math-frontier} skill: freeze the
claim, identify its first unresolved dependency, and specify how a route
could fail. Conceptual prompts requested a small number of distinct
mechanisms; paper-only tasks sought complete arguments before implementation.
A condensed representative prompt, adapted from the archived protocols, is:
\begin{quote}\small
Work on one precisely stated unresolved claim. Propose at most three
distinct mechanisms, each with an intermediate lemma, a noncircular
justification, and an adversarial test. Preserve the quantifiers and
right-sided cofactors. For a paper task, supply a complete argument; for
formalization, prove the required hypotheses without adding axioms or
using \texttt{sorry}. Return evidence and the first unsupported bridge; do not
promote a conditional result to a theorem.
\end{quote}

Recorded research and review sessions used GPT-5.5, GPT-5.6 Luna and Sol,
GPT-6 Astra, and Claude Fable 5, Haiku 4.5 and Opus 5. The project protocols
assigned routine implementation and replay to Luna and conceptual
consultation to Sol; independent reports also record Opus reviews.
Subsequent manuscript, proof-map, build and release work used Fable 5.1,
Opus 5, Sonnet 5 and GPT-6 Luna. These records establish participation and task roles,
without assigning a particular mathematical idea to one model.

The completed proof uses Lean~4 and Mathlib and pins AlgebraicAnalysis v0.3.0
\cite{Lean4,Mathlib2020,AlgebraicAnalysis030}; the later library archive is v0.3.2~\cite{AlgebraicAnalysis032}.
Formalization distinguished theorems proving required inputs from
conditional theorems merely consuming them. Statement comparison and
independent NanoDa and Lean kernel checks verified the submitted declarations
\cite{Comparator,NanoDa,PalomarDocs}; their proofs use only
\texttt{propext}, \texttt{Classical.choice}, and \texttt{Quot.sound}.
The Stafford38 and Global Stafford headline proofs are registered separately in
Palomar~\cite{Stafford38Palomar2026,GlobalStaffordPalomar2026}; their Lean developments are archived~\cite{Stafford38Zenodo120,GlobalStaffordZenodo106}.
These checks certify the formal declarations; human review of the paper's
exposition and its correspondence with them remains a separate task.

\else
\paragraph{AI-assisted research and formalization.}
Under Christopher Albert's direction, AI agents supported discovery,
counterexample searches, independent review, and formalization. Early
cofactor-selection approaches stalled; the proof advanced by reframing
the problem as vanishing of the canonical quotient, using Euler
surjectivity, monicity and characteristic-support geometry. Full Lean
formalization followed the completed paper argument, while bounded checks
continued during discovery. Parallel workers returned scoped evidence to
one controller, and reviewers checked quantifiers, sidedness and unsupported
bridges independently. The Lean/Mathlib development uses AlgebraicAnalysis
\cite{Lean4,Mathlib2020,AlgebraicAnalysis030}; its registered proof and
statement/kernel checks are recorded by Palomar
\cite{Stafford38Palomar2026,Comparator,NanoDa,PalomarDocs}.
The appendix describes the models, proof history, and prompt strategy.
\fi
\endgroup
% The manuscript contains an inline adaptation for direct editing.

\ifdefined\StaffordWorkflowAppendix
% Evidence: project ledgers, archived outward session messages, model-use
% records, independent reviews, and canonical formal-proof provenance.
\section*{Appendix: AI-assisted research and formalization}

Under Christopher Albert's direction, AI agents supported proof discovery,
counterexample searches, review, and formalization. The work began with
cyclic-extension computations in Macaulay2 and Maple. Attempts to construct
universal cofactors directly produced special cases and reductions but
eventually stalled. The decisive change was to study the canonical quotient
$Q=A_n/(dA_n+x^N dA_n)$: Euler surjectivity and monicity constrained its
characteristic support, while Gabber's involutivity and asymptotic conormal
geometry forced that support to be empty. This turned cofactor selection
into a quotient-vanishing problem. Full Lean formalization was deferred
until the paper argument was complete, although small formal checks
continued throughout. In the final formal development, support avoidance
over the coordinate hyperplane was proved by a localized filtered Koszul
length contradiction.

The workflow separated proposed ideas from verified claims. One controller
maintained the authoritative proof state, while parallel workers explored
distinct mechanisms and returned evidence with an explicit scope.
Independent reviewers checked fixed revisions for quantifiers, sidedness,
circularity, and missing hypotheses. The project protocols followed the
same discipline as the shared \texttt{math-frontier} skill: freeze the
claim, identify its first unresolved dependency, and specify how a route
could fail. Conceptual prompts requested a small number of distinct
mechanisms; paper-only tasks sought complete arguments before implementation.
A condensed representative prompt, adapted from the archived protocols, is:
\begin{quote}\small
Work on one precisely stated unresolved claim. Propose at most three
distinct mechanisms, each with an intermediate lemma, a noncircular
justification, and an adversarial test. Preserve the quantifiers and
right-sided cofactors. For a paper task, supply a complete argument; for
formalization, prove the required hypotheses without adding axioms or
using \texttt{sorry}. Return evidence and the first unsupported bridge; do not
promote a conditional result to a theorem.
\end{quote}

Recorded research and review sessions used GPT-5.5, GPT-5.6 Luna and Sol,
GPT-6 Astra, and Claude Fable 5, Haiku 4.5 and Opus 5. The project protocols
assigned routine implementation and replay to Luna and conceptual
consultation to Sol; independent reports also record Opus reviews.
Subsequent manuscript, proof-map, build and release work used Fable 5.1,
Opus 5, Sonnet 5 and GPT-6 Luna. These records establish participation and task roles,
without assigning a particular mathematical idea to one model.

The completed proof uses Lean~4 and Mathlib and pins AlgebraicAnalysis v0.3.0
\cite{Lean4,Mathlib2020,AlgebraicAnalysis030}; the later library archive is v0.3.2~\cite{AlgebraicAnalysis032}.
Formalization distinguished theorems proving required inputs from
conditional theorems merely consuming them. Statement comparison and
independent NanoDa and Lean kernel checks verified the submitted declarations
\cite{Comparator,NanoDa,PalomarDocs}; their proofs use only
\texttt{propext}, \texttt{Classical.choice}, and \texttt{Quot.sound}.
The Stafford38 and Global Stafford headline proofs are registered separately in
Palomar~\cite{Stafford38Palomar2026,GlobalStaffordPalomar2026}; their Lean developments are archived~\cite{Stafford38Zenodo120,GlobalStaffordZenodo106}.
These checks certify the formal declarations; human review of the paper's
exposition and its correspondence with them remains a separate task.

\else
\paragraph{AI-assisted research and formalization.}
Under Christopher Albert's direction, AI agents supported discovery,
counterexample searches, independent review, and formalization. Early
cofactor-selection approaches stalled; the proof advanced by reframing
the problem as vanishing of the canonical quotient, using Euler
surjectivity, monicity and characteristic-support geometry. Full Lean
formalization followed the completed paper argument, while bounded checks
continued during discovery. Parallel workers returned scoped evidence to
one controller, and reviewers checked quantifiers, sidedness and unsupported
bridges independently. The Lean/Mathlib development uses AlgebraicAnalysis
\cite{Lean4,Mathlib2020,AlgebraicAnalysis030}; its registered proof and
statement/kernel checks are recorded by Palomar
\cite{Stafford38Palomar2026,Comparator,NanoDa,PalomarDocs}.
The appendix describes the models, proof history, and prompt strategy.
\fi

% Optional appendix: \begingroup\def\StaffordWorkflowAppendix{}
% % Optional short paragraph: % Optional short paragraph: \input{ai_workflow}
% Optional appendix: \begingroup\def\StaffordWorkflowAppendix{}
% \input{ai_workflow}\endgroup
% The manuscript contains an inline adaptation for direct editing.

\ifdefined\StaffordWorkflowAppendix
% Evidence: project ledgers, archived outward session messages, model-use
% records, independent reviews, and canonical formal-proof provenance.
\section*{Appendix: AI-assisted research and formalization}

Under Christopher Albert's direction, AI agents supported proof discovery,
counterexample searches, review, and formalization. The work began with
cyclic-extension computations in Macaulay2 and Maple. Attempts to construct
universal cofactors directly produced special cases and reductions but
eventually stalled. The decisive change was to study the canonical quotient
$Q=A_n/(dA_n+x^N dA_n)$: Euler surjectivity and monicity constrained its
characteristic support, while Gabber's involutivity and asymptotic conormal
geometry forced that support to be empty. This turned cofactor selection
into a quotient-vanishing problem. Full Lean formalization was deferred
until the paper argument was complete, although small formal checks
continued throughout. In the final formal development, support avoidance
over the coordinate hyperplane was proved by a localized filtered Koszul
length contradiction.

The workflow separated proposed ideas from verified claims. One controller
maintained the authoritative proof state, while parallel workers explored
distinct mechanisms and returned evidence with an explicit scope.
Independent reviewers checked fixed revisions for quantifiers, sidedness,
circularity, and missing hypotheses. The project protocols followed the
same discipline as the shared \texttt{math-frontier} skill: freeze the
claim, identify its first unresolved dependency, and specify how a route
could fail. Conceptual prompts requested a small number of distinct
mechanisms; paper-only tasks sought complete arguments before implementation.
A condensed representative prompt, adapted from the archived protocols, is:
\begin{quote}\small
Work on one precisely stated unresolved claim. Propose at most three
distinct mechanisms, each with an intermediate lemma, a noncircular
justification, and an adversarial test. Preserve the quantifiers and
right-sided cofactors. For a paper task, supply a complete argument; for
formalization, prove the required hypotheses without adding axioms or
using \texttt{sorry}. Return evidence and the first unsupported bridge; do not
promote a conditional result to a theorem.
\end{quote}

Recorded research and review sessions used GPT-5.5, GPT-5.6 Luna and Sol,
GPT-6 Astra, and Claude Fable 5, Haiku 4.5 and Opus 5. The project protocols
assigned routine implementation and replay to Luna and conceptual
consultation to Sol; independent reports also record Opus reviews.
Subsequent manuscript, proof-map, build and release work used Fable 5.1,
Opus 5, Sonnet 5 and GPT-6 Luna. These records establish participation and task roles,
without assigning a particular mathematical idea to one model.

The completed proof uses Lean~4 and Mathlib and pins AlgebraicAnalysis v0.3.0
\cite{Lean4,Mathlib2020,AlgebraicAnalysis030}; the later library archive is v0.3.2~\cite{AlgebraicAnalysis032}.
Formalization distinguished theorems proving required inputs from
conditional theorems merely consuming them. Statement comparison and
independent NanoDa and Lean kernel checks verified the submitted declarations
\cite{Comparator,NanoDa,PalomarDocs}; their proofs use only
\texttt{propext}, \texttt{Classical.choice}, and \texttt{Quot.sound}.
The Stafford38 and Global Stafford headline proofs are registered separately in
Palomar~\cite{Stafford38Palomar2026,GlobalStaffordPalomar2026}; their Lean developments are archived~\cite{Stafford38Zenodo120,GlobalStaffordZenodo106}.
These checks certify the formal declarations; human review of the paper's
exposition and its correspondence with them remains a separate task.

\else
\paragraph{AI-assisted research and formalization.}
Under Christopher Albert's direction, AI agents supported discovery,
counterexample searches, independent review, and formalization. Early
cofactor-selection approaches stalled; the proof advanced by reframing
the problem as vanishing of the canonical quotient, using Euler
surjectivity, monicity and characteristic-support geometry. Full Lean
formalization followed the completed paper argument, while bounded checks
continued during discovery. Parallel workers returned scoped evidence to
one controller, and reviewers checked quantifiers, sidedness and unsupported
bridges independently. The Lean/Mathlib development uses AlgebraicAnalysis
\cite{Lean4,Mathlib2020,AlgebraicAnalysis030}; its registered proof and
statement/kernel checks are recorded by Palomar
\cite{Stafford38Palomar2026,Comparator,NanoDa,PalomarDocs}.
The appendix describes the models, proof history, and prompt strategy.
\fi

% Optional appendix: \begingroup\def\StaffordWorkflowAppendix{}
% % Optional short paragraph: \input{ai_workflow}
% Optional appendix: \begingroup\def\StaffordWorkflowAppendix{}
% \input{ai_workflow}\endgroup
% The manuscript contains an inline adaptation for direct editing.

\ifdefined\StaffordWorkflowAppendix
% Evidence: project ledgers, archived outward session messages, model-use
% records, independent reviews, and canonical formal-proof provenance.
\section*{Appendix: AI-assisted research and formalization}

Under Christopher Albert's direction, AI agents supported proof discovery,
counterexample searches, review, and formalization. The work began with
cyclic-extension computations in Macaulay2 and Maple. Attempts to construct
universal cofactors directly produced special cases and reductions but
eventually stalled. The decisive change was to study the canonical quotient
$Q=A_n/(dA_n+x^N dA_n)$: Euler surjectivity and monicity constrained its
characteristic support, while Gabber's involutivity and asymptotic conormal
geometry forced that support to be empty. This turned cofactor selection
into a quotient-vanishing problem. Full Lean formalization was deferred
until the paper argument was complete, although small formal checks
continued throughout. In the final formal development, support avoidance
over the coordinate hyperplane was proved by a localized filtered Koszul
length contradiction.

The workflow separated proposed ideas from verified claims. One controller
maintained the authoritative proof state, while parallel workers explored
distinct mechanisms and returned evidence with an explicit scope.
Independent reviewers checked fixed revisions for quantifiers, sidedness,
circularity, and missing hypotheses. The project protocols followed the
same discipline as the shared \texttt{math-frontier} skill: freeze the
claim, identify its first unresolved dependency, and specify how a route
could fail. Conceptual prompts requested a small number of distinct
mechanisms; paper-only tasks sought complete arguments before implementation.
A condensed representative prompt, adapted from the archived protocols, is:
\begin{quote}\small
Work on one precisely stated unresolved claim. Propose at most three
distinct mechanisms, each with an intermediate lemma, a noncircular
justification, and an adversarial test. Preserve the quantifiers and
right-sided cofactors. For a paper task, supply a complete argument; for
formalization, prove the required hypotheses without adding axioms or
using \texttt{sorry}. Return evidence and the first unsupported bridge; do not
promote a conditional result to a theorem.
\end{quote}

Recorded research and review sessions used GPT-5.5, GPT-5.6 Luna and Sol,
GPT-6 Astra, and Claude Fable 5, Haiku 4.5 and Opus 5. The project protocols
assigned routine implementation and replay to Luna and conceptual
consultation to Sol; independent reports also record Opus reviews.
Subsequent manuscript, proof-map, build and release work used Fable 5.1,
Opus 5, Sonnet 5 and GPT-6 Luna. These records establish participation and task roles,
without assigning a particular mathematical idea to one model.

The completed proof uses Lean~4 and Mathlib and pins AlgebraicAnalysis v0.3.0
\cite{Lean4,Mathlib2020,AlgebraicAnalysis030}; the later library archive is v0.3.2~\cite{AlgebraicAnalysis032}.
Formalization distinguished theorems proving required inputs from
conditional theorems merely consuming them. Statement comparison and
independent NanoDa and Lean kernel checks verified the submitted declarations
\cite{Comparator,NanoDa,PalomarDocs}; their proofs use only
\texttt{propext}, \texttt{Classical.choice}, and \texttt{Quot.sound}.
The Stafford38 and Global Stafford headline proofs are registered separately in
Palomar~\cite{Stafford38Palomar2026,GlobalStaffordPalomar2026}; their Lean developments are archived~\cite{Stafford38Zenodo120,GlobalStaffordZenodo106}.
These checks certify the formal declarations; human review of the paper's
exposition and its correspondence with them remains a separate task.

\else
\paragraph{AI-assisted research and formalization.}
Under Christopher Albert's direction, AI agents supported discovery,
counterexample searches, independent review, and formalization. Early
cofactor-selection approaches stalled; the proof advanced by reframing
the problem as vanishing of the canonical quotient, using Euler
surjectivity, monicity and characteristic-support geometry. Full Lean
formalization followed the completed paper argument, while bounded checks
continued during discovery. Parallel workers returned scoped evidence to
one controller, and reviewers checked quantifiers, sidedness and unsupported
bridges independently. The Lean/Mathlib development uses AlgebraicAnalysis
\cite{Lean4,Mathlib2020,AlgebraicAnalysis030}; its registered proof and
statement/kernel checks are recorded by Palomar
\cite{Stafford38Palomar2026,Comparator,NanoDa,PalomarDocs}.
The appendix describes the models, proof history, and prompt strategy.
\fi
\endgroup
% The manuscript contains an inline adaptation for direct editing.

\ifdefined\StaffordWorkflowAppendix
% Evidence: project ledgers, archived outward session messages, model-use
% records, independent reviews, and canonical formal-proof provenance.
\section*{Appendix: AI-assisted research and formalization}

Under Christopher Albert's direction, AI agents supported proof discovery,
counterexample searches, review, and formalization. The work began with
cyclic-extension computations in Macaulay2 and Maple. Attempts to construct
universal cofactors directly produced special cases and reductions but
eventually stalled. The decisive change was to study the canonical quotient
$Q=A_n/(dA_n+x^N dA_n)$: Euler surjectivity and monicity constrained its
characteristic support, while Gabber's involutivity and asymptotic conormal
geometry forced that support to be empty. This turned cofactor selection
into a quotient-vanishing problem. Full Lean formalization was deferred
until the paper argument was complete, although small formal checks
continued throughout. In the final formal development, support avoidance
over the coordinate hyperplane was proved by a localized filtered Koszul
length contradiction.

The workflow separated proposed ideas from verified claims. One controller
maintained the authoritative proof state, while parallel workers explored
distinct mechanisms and returned evidence with an explicit scope.
Independent reviewers checked fixed revisions for quantifiers, sidedness,
circularity, and missing hypotheses. The project protocols followed the
same discipline as the shared \texttt{math-frontier} skill: freeze the
claim, identify its first unresolved dependency, and specify how a route
could fail. Conceptual prompts requested a small number of distinct
mechanisms; paper-only tasks sought complete arguments before implementation.
A condensed representative prompt, adapted from the archived protocols, is:
\begin{quote}\small
Work on one precisely stated unresolved claim. Propose at most three
distinct mechanisms, each with an intermediate lemma, a noncircular
justification, and an adversarial test. Preserve the quantifiers and
right-sided cofactors. For a paper task, supply a complete argument; for
formalization, prove the required hypotheses without adding axioms or
using \texttt{sorry}. Return evidence and the first unsupported bridge; do not
promote a conditional result to a theorem.
\end{quote}

Recorded research and review sessions used GPT-5.5, GPT-5.6 Luna and Sol,
GPT-6 Astra, and Claude Fable 5, Haiku 4.5 and Opus 5. The project protocols
assigned routine implementation and replay to Luna and conceptual
consultation to Sol; independent reports also record Opus reviews.
Subsequent manuscript, proof-map, build and release work used Fable 5.1,
Opus 5, Sonnet 5 and GPT-6 Luna. These records establish participation and task roles,
without assigning a particular mathematical idea to one model.

The completed proof uses Lean~4 and Mathlib and pins AlgebraicAnalysis v0.3.0
\cite{Lean4,Mathlib2020,AlgebraicAnalysis030}; the later library archive is v0.3.2~\cite{AlgebraicAnalysis032}.
Formalization distinguished theorems proving required inputs from
conditional theorems merely consuming them. Statement comparison and
independent NanoDa and Lean kernel checks verified the submitted declarations
\cite{Comparator,NanoDa,PalomarDocs}; their proofs use only
\texttt{propext}, \texttt{Classical.choice}, and \texttt{Quot.sound}.
The Stafford38 and Global Stafford headline proofs are registered separately in
Palomar~\cite{Stafford38Palomar2026,GlobalStaffordPalomar2026}; their Lean developments are archived~\cite{Stafford38Zenodo120,GlobalStaffordZenodo106}.
These checks certify the formal declarations; human review of the paper's
exposition and its correspondence with them remains a separate task.

\else
\paragraph{AI-assisted research and formalization.}
Under Christopher Albert's direction, AI agents supported discovery,
counterexample searches, independent review, and formalization. Early
cofactor-selection approaches stalled; the proof advanced by reframing
the problem as vanishing of the canonical quotient, using Euler
surjectivity, monicity and characteristic-support geometry. Full Lean
formalization followed the completed paper argument, while bounded checks
continued during discovery. Parallel workers returned scoped evidence to
one controller, and reviewers checked quantifiers, sidedness and unsupported
bridges independently. The Lean/Mathlib development uses AlgebraicAnalysis
\cite{Lean4,Mathlib2020,AlgebraicAnalysis030}; its registered proof and
statement/kernel checks are recorded by Palomar
\cite{Stafford38Palomar2026,Comparator,NanoDa,PalomarDocs}.
The appendix describes the models, proof history, and prompt strategy.
\fi
\endgroup
% The manuscript contains an inline adaptation for direct editing.

\ifdefined\StaffordWorkflowAppendix
% Evidence: project ledgers, archived outward session messages, model-use
% records, independent reviews, and canonical formal-proof provenance.
\section*{Appendix: AI-assisted research and formalization}

Under Christopher Albert's direction, AI agents supported proof discovery,
counterexample searches, review, and formalization. The work began with
cyclic-extension computations in Macaulay2 and Maple. Attempts to construct
universal cofactors directly produced special cases and reductions but
eventually stalled. The decisive change was to study the canonical quotient
$Q=A_n/(dA_n+x^N dA_n)$: Euler surjectivity and monicity constrained its
characteristic support, while Gabber's involutivity and asymptotic conormal
geometry forced that support to be empty. This turned cofactor selection
into a quotient-vanishing problem. Full Lean formalization was deferred
until the paper argument was complete, although small formal checks
continued throughout. In the final formal development, support avoidance
over the coordinate hyperplane was proved by a localized filtered Koszul
length contradiction.

The workflow separated proposed ideas from verified claims. One controller
maintained the authoritative proof state, while parallel workers explored
distinct mechanisms and returned evidence with an explicit scope.
Independent reviewers checked fixed revisions for quantifiers, sidedness,
circularity, and missing hypotheses. The project protocols followed the
same discipline as the shared \texttt{math-frontier} skill: freeze the
claim, identify its first unresolved dependency, and specify how a route
could fail. Conceptual prompts requested a small number of distinct
mechanisms; paper-only tasks sought complete arguments before implementation.
A condensed representative prompt, adapted from the archived protocols, is:
\begin{quote}\small
Work on one precisely stated unresolved claim. Propose at most three
distinct mechanisms, each with an intermediate lemma, a noncircular
justification, and an adversarial test. Preserve the quantifiers and
right-sided cofactors. For a paper task, supply a complete argument; for
formalization, prove the required hypotheses without adding axioms or
using \texttt{sorry}. Return evidence and the first unsupported bridge; do not
promote a conditional result to a theorem.
\end{quote}

Recorded research and review sessions used GPT-5.5, GPT-5.6 Luna and Sol,
GPT-6 Astra, and Claude Fable 5, Haiku 4.5 and Opus 5. The project protocols
assigned routine implementation and replay to Luna and conceptual
consultation to Sol; independent reports also record Opus reviews.
Subsequent manuscript, proof-map, build and release work used Fable 5.1,
Opus 5, Sonnet 5 and GPT-6 Luna. These records establish participation and task roles,
without assigning a particular mathematical idea to one model.

The completed proof uses Lean~4 and Mathlib and pins AlgebraicAnalysis v0.3.0
\cite{Lean4,Mathlib2020,AlgebraicAnalysis030}; the later library archive is v0.3.2~\cite{AlgebraicAnalysis032}.
Formalization distinguished theorems proving required inputs from
conditional theorems merely consuming them. Statement comparison and
independent NanoDa and Lean kernel checks verified the submitted declarations
\cite{Comparator,NanoDa,PalomarDocs}; their proofs use only
\texttt{propext}, \texttt{Classical.choice}, and \texttt{Quot.sound}.
The Stafford38 and Global Stafford headline proofs are registered separately in
Palomar~\cite{Stafford38Palomar2026,GlobalStaffordPalomar2026}; their Lean developments are archived~\cite{Stafford38Zenodo120,GlobalStaffordZenodo106}.
These checks certify the formal declarations; human review of the paper's
exposition and its correspondence with them remains a separate task.

\else
\paragraph{AI-assisted research and formalization.}
Under Christopher Albert's direction, AI agents supported discovery,
counterexample searches, independent review, and formalization. Early
cofactor-selection approaches stalled; the proof advanced by reframing
the problem as vanishing of the canonical quotient, using Euler
surjectivity, monicity and characteristic-support geometry. Full Lean
formalization followed the completed paper argument, while bounded checks
continued during discovery. Parallel workers returned scoped evidence to
one controller, and reviewers checked quantifiers, sidedness and unsupported
bridges independently. The Lean/Mathlib development uses AlgebraicAnalysis
\cite{Lean4,Mathlib2020,AlgebraicAnalysis030}; its registered proof and
statement/kernel checks are recorded by Palomar
\cite{Stafford38Palomar2026,Comparator,NanoDa,PalomarDocs}.
The appendix describes the models, proof history, and prompt strategy.
\fi
\endgroup
% The manuscript contains an inline adaptation for direct editing.

\ifdefined\StaffordWorkflowAppendix
% Evidence: project ledgers, archived outward session messages, model-use
% records, independent reviews, and canonical formal-proof provenance.
\section*{Appendix: AI-assisted research and formalization}

Under Christopher Albert's direction, AI agents supported proof discovery,
counterexample searches, review, and formalization. The work began with
cyclic-extension computations in Macaulay2 and Maple. Attempts to construct
universal cofactors directly produced special cases and reductions but
eventually stalled. The decisive change was to study the canonical quotient
$Q=A_n/(dA_n+x^N dA_n)$: Euler surjectivity and monicity constrained its
characteristic support, while Gabber's involutivity and asymptotic conormal
geometry forced that support to be empty. This turned cofactor selection
into a quotient-vanishing problem. Full Lean formalization was deferred
until the paper argument was complete, although small formal checks
continued throughout. In the final formal development, support avoidance
over the coordinate hyperplane was proved by a localized filtered Koszul
length contradiction.

The workflow separated proposed ideas from verified claims. One controller
maintained the authoritative proof state, while parallel workers explored
distinct mechanisms and returned evidence with an explicit scope.
Independent reviewers checked fixed revisions for quantifiers, sidedness,
circularity, and missing hypotheses. The project protocols followed the
same discipline as the shared \texttt{math-frontier} skill: freeze the
claim, identify its first unresolved dependency, and specify how a route
could fail. Conceptual prompts requested a small number of distinct
mechanisms; paper-only tasks sought complete arguments before implementation.
A condensed representative prompt, adapted from the archived protocols, is:
\begin{quote}\small
Work on one precisely stated unresolved claim. Propose at most three
distinct mechanisms, each with an intermediate lemma, a noncircular
justification, and an adversarial test. Preserve the quantifiers and
right-sided cofactors. For a paper task, supply a complete argument; for
formalization, prove the required hypotheses without adding axioms or
using \texttt{sorry}. Return evidence and the first unsupported bridge; do not
promote a conditional result to a theorem.
\end{quote}

Recorded research and review sessions used GPT-5.5, GPT-5.6 Luna and Sol,
GPT-6 Astra, and Claude Fable 5, Haiku 4.5 and Opus 5. The project protocols
assigned routine implementation and replay to Luna and conceptual
consultation to Sol; independent reports also record Opus reviews.
Subsequent manuscript, proof-map, build and release work used Fable 5.1,
Opus 5, Sonnet 5 and GPT-6 Luna. These records establish participation and task roles,
without assigning a particular mathematical idea to one model.

The completed proof uses Lean~4 and Mathlib and pins AlgebraicAnalysis v0.3.0
\cite{Lean4,Mathlib2020,AlgebraicAnalysis030}; the later library archive is v0.3.2~\cite{AlgebraicAnalysis032}.
Formalization distinguished theorems proving required inputs from
conditional theorems merely consuming them. Statement comparison and
independent NanoDa and Lean kernel checks verified the submitted declarations
\cite{Comparator,NanoDa,PalomarDocs}; their proofs use only
\texttt{propext}, \texttt{Classical.choice}, and \texttt{Quot.sound}.
The Stafford38 and Global Stafford headline proofs are registered separately in
Palomar~\cite{Stafford38Palomar2026,GlobalStaffordPalomar2026}; their Lean developments are archived~\cite{Stafford38Zenodo120,GlobalStaffordZenodo106}.
These checks certify the formal declarations; human review of the paper's
exposition and its correspondence with them remains a separate task.

\else
\paragraph{AI-assisted research and formalization.}
Under Christopher Albert's direction, AI agents supported discovery,
counterexample searches, independent review, and formalization. Early
cofactor-selection approaches stalled; the proof advanced by reframing
the problem as vanishing of the canonical quotient, using Euler
surjectivity, monicity and characteristic-support geometry. Full Lean
formalization followed the completed paper argument, while bounded checks
continued during discovery. Parallel workers returned scoped evidence to
one controller, and reviewers checked quantifiers, sidedness and unsupported
bridges independently. The Lean/Mathlib development uses AlgebraicAnalysis
\cite{Lean4,Mathlib2020,AlgebraicAnalysis030}; its registered proof and
statement/kernel checks are recorded by Palomar
\cite{Stafford38Palomar2026,Comparator,NanoDa,PalomarDocs}.
The appendix describes the models, proof history, and prompt strategy.
\fi
